\begin{textsample}{POS Dim 3 – to – Score 17.00 – to/to\_inf\_25.txt}  \label{ex:f3_pos_097}
A documentação pedagógica consiste na observação e \textbf{registro} do que fazem as \textbf{crianças} da Educação \textbf{Infantil}.
A prática é diversificada, porém, não é plural, ou seja, há muitas formas de realizá -la, mas apenas dois focos estruturantes : tornar visível e evidente as aprendizagens e a quem elas se destinam.
Documentação pedagógica É verdade que muitos de nossos professores já realizam algum tipo de \textbf{registro}, mas sua simples realização não se converte em uma verdadeira documentação pedagógica, uma vez que ela, além de dar visibilidade aos fazeres \textbf{infantis}, “ possibilita realizar, analisar e problematizar, de forma pública ou coletiva, aquilo que foi observado e registrado, assim como a inseparabilidade entre o \textbf{documentado} e o processo de planejamento, a definição do currículo, a escolha das atividades, a participação das \textbf{crianças} e das famílias no processo de documentação ” ( BARBOSA ; FERNANDES, 2012 ).
O primeiro passo para a elaboração da documentação pedagógica é a elaboração de algum tipo de \textbf{registro} ( anotações, fotos, filmes, gravações e produções das \textbf{crianças} ), que conterá as informações levantadas nas observações.
Alguns questionamentos precisam ser realizados pelos professores : O que dizem essas informações?
O que se descobriu com essas evidências?
Que informações respondem às inquietações docentes?
O que as \textbf{crianças} realmente aprenderam?
Essas reflexões são o primeiro passo da documentação.
Em seguida, é preciso ter claro a quem se quer comunicar essas informações.
Aí está a alma da documentação pedagógica, ou seja, a comunicação.
Os professores precisam ter claro quem é o público-alvo, para assim, selecionar os conteúdos dos \textbf{registros} e as formas de apresentá -los.
Há vários públicos, com dimensões e intenções distintas, que vão desde o próprio educador, a equipe pedagógica, as \textbf{crianças} com as quais atuam, suas respectivas famílias e a comunidade em geral.
As formas de \textbf{documentar} os processos de aprendizagens, e das vivências cotidianas das \textbf{crianças}, sugerem uma nova forma de conceber o planejamento, o currículo e a avaliação.
Nessa perspectiva, principalmente, ao que tange a avaliação, ela só tem sentido na análise dos componentes do processo, não há nenhuma intenção de categorizar os \textbf{pequenos}, de acordo com algum nível de desenvolvimento.
O foco está no processo das suas interações e aprendizagens.
O eu, o outro e o nós SÍNTESE DAS APRENDIZAGENS - respeitar e expressar sentimentos e emoções ; - atuar em grupo e \textbf{demonstrar} interesse em construir novas relações, respeitando a diversidade e solidarizando -se com os outros ; - conhecer e respeitar regras de convívio social, manifestando respeito pelo outro. - reconhecer a importância de ações e situações do cotidiano que contribuem para o \textbf{cuidado} de sua saúde e a manutenção de ambientes saudáveis ; - apresentar autonomia nas práticas de \textbf{higiene}, alimentação, vestir -se e no \textbf{cuidado} com seu bem-estar, valorizando o próprio corpo ; - utilizar o corpo intencionalmente ( com criatividade, controle e adequação ) como instrumento de interação com o outro e com o meio ; - coordenar suas habilidades manuais.
Traços, \textbf{sons}, cores e formas - discriminar os diferentes tipos de \textbf{sons} e ritmos e interagir com a música, percebendo -a como forma de expressão individual e coletiva ; - expressar -se por meio das artes visuais, utilizando diferentes materiais ; - relacionar -se com o outro empregando \textbf{gestos}, palavras, \textbf{brincadeiras}, jogos, \textbf{imitações}, observações e expressão corporal. \textbf{Escuta}, fala, pensamento e imaginação - Expressar ideias, \textbf{desejos} e sentimentos em distintas situações de interação, por diferentes meios ; - Argumentar e relatar fatos oralmente, em sequência temporal e causal, organizando e adequando sua fala ao contexto em que é produzida ; - Ouvir, compreender, \textbf{contar}, recontar e criar narrativas ; - Conhecer diferentes gêneros e portadores textuais, \textbf{demonstrando} compreensão da função social da escrita e reconhecendo a leitura como fonte de prazer e informação.
Espaços, tempos, quantidades, relações e transformações - identificar, nomear adequadamente e comparar as propriedades dos objetos, estabelecendo relações entre eles ; - interagir com o meio ambiente e com fenômenos naturais ou artificiais, \textbf{demonstrando} \textbf{curiosidade} e \textbf{cuidado} com relação a eles ; - utilizar vocabulário relativo às noções de grandeza ( maior, \textbf{menor}, igual etc. ), espaço ( dentro e fora ) e medidas ( comprido, curto, grosso, fino ) como meio de comunicação de suas experiências ; - utilizar unidades de medida ( dia e noite ; dias, semanas, \textbf{meses} e ano ) e noções de tempo ( presente, passado e futuro ; antes, agora e depois ), para responder a necessidades e questões do cotidiano ; - identificar e registrar quantidades por meio de diferentes formas de representação ( contagens, desenhos, símbolos, escrita de números, organização de gráficos básicos etc. ). ( BRASIL, 2017, p. 52 e 53 )

% matched lemmas: brincadeira, contar, criança, cuidado, curiosidade, demonstrar, desejo, documentar, escuta, gesto, higiene, imitação, infantil, mês, pequeno, registro, som
\end{textsample}
